\chapter{Grille complète des résultats}
\label{ann:resultats}
% BUDGET : 2,5 pages. Version 2026-09-02 : l'annexe ne contient plus que des tableaux. Les
% reserves de lecture, le statut des fronts et les prix duals, jusqu'ici en prose, sont devenus
% trois tables.
%
% Regenerer d'abord :  .venv/Scripts/python memoire/build_chiffres.py
% Tous les nombres de cette annexe sont des macros lues dans les recap.json des runs, sauf les
% trois prix duals.
%
% NE PAS remettre \gls{hhi} dans une cellule de la table ci-dessous : c'est sa premiere
% occurrence du document, la forme longue s'y developpe et fait deborder le tableau de 122 pt.
% Elle est appelee depuis la LEGENDE, qui est un paragraphe et peut donc la couper.

\section{Les politiques non forcées, tous indicateurs}

\begin{table}[H]
  \centering
  \caption{Économie, indicateurs physiques et convergence des sept politiques qui ont abouti
  sans forçage. La concentration est mesurée par l'indice \gls{hhi}. La colonne P1 porte un
  autre objectif, la marge brute pure, si bien que sa ligne d'objectif n'est pas comparable aux
  autres et que ses lignes physiques le sont. L'eau ne porte que sur les parcelles irrigables
  et les \glsentryshort{ges} sont dans l'unité interne du modèle.}
  \label{tab:ann-pros-eco}
  \scriptsize
  % Huit colonnes de nombres a sept chiffres : le tabcolsep par defaut (6 pt, soit 84 pt
  % cumules) fait deborder le tableau de 36 pt. A 3 pt il tient sans reduire la fonte.
  \setlength{\tabcolsep}{3pt}
  \begin{tabular}{@{}lrrrrrrr@{}}
    \toprule
    & \textbf{P1} & \textbf{P4} & \textbf{P5} & \textbf{P6} & \textbf{P7} & \textbf{P8}
      & \textbf{P9} \\
    \midrule
    Objectif (M\eur) & \num{\prosPunObjectifM} & \num{\prosPquatreObjectifM}
      & \num{\prosPcinqObjectifM} & \num{\prosPsixObjectifM} & \num{\prosPseptObjectifM}
      & \num{\prosPhuitObjectifM} & \num{\prosPneufObjectifM} \\
    Marge brute (M\eur) & \num{\prosPunMargeM} & \num{\prosPquatreMargeM} & \num{\prosPcinqMargeM}
      & \num{\prosPsixMargeM} & \num{\prosPseptMargeM} & \num{\prosPhuitMargeM}
      & \num{\prosPneufMargeM} \\
    Subvention (M\eur) & \num{\prosPunSubventionM} & \num{\prosPquatreSubventionM}
      & \num{\prosPcinqSubventionM} & \num{\prosPsixSubventionM} & \num{\prosPseptSubventionM}
      & \num{\prosPhuitSubventionM} & \num{\prosPneufSubventionM} \\
    Emploi (\glsentryshort{etp}) & \num{\prosPunEtp} & \num{\prosPquatreEtp} & \num{\prosPcinqEtp}
      & \num{\prosPsixEtp} & \num{\prosPseptEtp} & \num{\prosPhuitEtp} & \num{\prosPneufEtp} \\
    Surface allouée (ha) & \num{\prosPunSurfaceHa} & \num{\prosPquatreSurfaceHa}
      & \num{\prosPcinqSurfaceHa} & \num{\prosPsixSurfaceHa} & \num{\prosPseptSurfaceHa}
      & \num{\prosPhuitSurfaceHa} & \num{\prosPneufSurfaceHa} \\
    Concentration (\glsentryshort{hhi}) & \num{\prosPunHhi} & \num{\prosPquatreHhi}
      & \num{\prosPcinqHhi} & \num{\prosPsixHhi} & \num{\prosPseptHhi} & \num{\prosPhuitHhi}
      & \num{\prosPneufHhi} \\
    \midrule
    Azote (kg N) & \num{\prosPunAzote} & \num{\prosPquatreAzote} & \num{\prosPcinqAzote}
      & \num{\prosPsixAzote} & \num{\prosPseptAzote} & \num{\prosPhuitAzote}
      & \num{\prosPneufAzote} \\
    \glsentryshort{ift} & \num{\prosPunIft} & \num{\prosPquatreIft} & \num{\prosPcinqIft}
      & \num{\prosPsixIft} & \num{\prosPseptIft} & \num{\prosPhuitIft} & \num{\prosPneufIft} \\
    \glsentryshort{ges} & \num{\prosPunGes} & \num{\prosPquatreGes} & \num{\prosPcinqGes}
      & \num{\prosPsixGes} & \num{\prosPseptGes} & \num{\prosPhuitGes} & \num{\prosPneufGes} \\
    Eau (m\textsuperscript{3}) & \num{\prosPunEau} & \num{\prosPquatreEau} & \num{\prosPcinqEau}
      & \num{\prosPsixEau} & \num{\prosPseptEau} & \num{\prosPhuitEau} & \num{\prosPneufEau} \\
    \midrule
    Durée (s) & \num{\prosPunDuree} & \num{\prosPquatreDuree} & \num{\prosPcinqDuree}
      & \num{\prosPsixDuree} & \num{\prosPseptDuree} & \num{\prosPhuitDuree}
      & \num{\prosPneufDuree} \\
    Optimum prouvé & \prosPunProuve & \prosPquatreProuve & \prosPcinqProuve & \prosPsixProuve
      & \prosPseptProuve & \prosPhuitProuve & \prosPneufProuve \\
    \bottomrule
  \end{tabular}
\end{table}

\begin{table}[H]
  \centering
  \caption{Les quatre réserves à appliquer à la lecture des tableaux de cette annexe.}
  \label{tab:ann-lectures}
  \scriptsize
  \setlength{\tabcolsep}{4pt}
  \begin{tabularx}{\textwidth}{@{}>{\raggedright\arraybackslash}p{34mm}X@{}}
    \toprule
    \textbf{Réserve} & \textbf{Contenu} \\
    \midrule
    Chaque cellule porte son drapeau de convergence & un solve arrêté à la limite de temps est
      prouvablement faux et non seulement imprécis (§~\ref{sec:plateau}). Un tableau de
      résultats qui ne le dit pas invite à comparer des chiffres qui ne sont pas comparables \\
    La surface allouée de P1 se lit avec sa marge & P1, qui supprime toute aide et toute règle,
      fait plus que doubler la marge brute du territoire tout en réduisant l'azote et la
      surface allouée. Elle concentre la production sur les meilleures terres et abandonne le
      reste, ce qu'aucun indicateur agrégé ne dirait \\
    Les durées ne se comparent pas d'une politique à l'autre & certains runs sont amorcés à
      chaud, et le même problème avec la même graine a demandé \SI{\mesHuitPreuveDuree}{\second}
      puis \SI{\mesSolveProsMax}{\second} \\
    Un indicateur qu'une contrainte fixe n'est pas un résultat & le plafond d'azote de P8 vaut
      exactement l'azote de P8, et le noter reviendrait à féliciter la politique d'avoir
      respecté sa propre consigne. Ces indicateurs sont détectés automatiquement depuis la
      liste des contraintes du run, donc rétroactivement sur les runs déjà écrits, et exclus de
      tout score composite \\
    \bottomrule
  \end{tabularx}
\end{table}

\section{Robustesse}

\begin{table}[H]
  \centering
  \caption{Regret par politique et par forçage, en millions d'euros d'écart à la meilleure
  politique du même forçage et jamais à un optimum global. Brut inclut la subvention, hors aide
  la déduit. Un tiret marque une cellule non résolue et non un résultat nul. La dernière ligne
  nomme la politique de référence de chaque colonne.}
  \label{tab:ann-regret}
  \small
  \begin{tabularx}{\textwidth}{@{}>{\raggedright\arraybackslash}Xrrrrrr@{}}
    \toprule
    & \multicolumn{2}{c}{\textbf{F0 nominal}} & \multicolumn{2}{c}{\textbf{F6 intrants}}
    & \multicolumn{2}{c}{\textbf{F9 crise}} \\
    \cmidrule(lr){2-3}\cmidrule(lr){4-5}\cmidrule(lr){6-7}
    & brut & hors aide & brut & hors aide & brut & hors aide \\
    \midrule
    P4 statu quo    & \num{\regretPquatreFzero} & \num{\regretNetPquatreFzero}
                    & \num{\regretPquatreFsix}  & \num{\regretNetPquatreFsix}
                    & \num{\regretPquatreFneuf} & \num{\regretNetPquatreFneuf} \\
    P6 verdissement & \num{\regretPsixFzero} & \num{\regretNetPsixFzero}
                    & \num{\regretPsixFsix}  & \num{\regretNetPsixFsix} & --- & --- \\
    P7 Écophyto     & \num{\regretPseptFzero} & \num{\regretNetPseptFzero}
                    & \num{\regretPseptFsix}  & \num{\regretNetPseptFsix} & --- & --- \\
    P8 agroécologie & \num{\regretPhuitFzero} & \num{\regretNetPhuitFzero} & --- & ---
                    & \num{\regretPhuitFneuf} & \num{\regretNetPhuitFneuf} \\
    P9 souveraineté & \num{\regretPneufFzero} & \num{\regretNetPneufFzero} & --- & ---
                    & \num{\regretPneufFneuf} & \num{\regretNetPneufFneuf} \\
    \midrule
    meilleure politique & \grilleMeilleureFzero & \grilleMeilleureNetFzero
      & \grilleMeilleureFsix & \grilleMeilleureNetFsix
      & \grilleMeilleureFneuf & \grilleMeilleureNetFneuf \\
    \bottomrule
  \end{tabularx}
\end{table}

\section{Les fronts, point par point}

\begin{table}[H]
  \centering
  \caption{Nombre de points de chaque front, nombre de points prouvés optimaux et nombre de
  paliers. Seul le front sur l'enveloppe publique est citable point par point.}
  \label{tab:ann-fronts-statut}
  \small
  \begin{tabularx}{\textwidth}{@{}>{\raggedright\arraybackslash}Xrrl@{}}
    \toprule
    \textbf{Front} & \textbf{Points} & \textbf{Prouvés} & \textbf{Paliers} \\
    \midrule
    Enveloppe publique, sous P4 & \num{\frontsubvQuatrePoints}
      & \num{\frontsubvQuatreProuves}
      & \num{\frontsubvQuatrePlateau} points de contrôle, plafond non contraignant \\
    Azote, configuration de référence & \num{\frontazoteRefPoints}
      & \num{\frontazoteRefProuves} & \num{\frontazoteRefPlateau} \\
    Azote, sous P8 & \num{\frontazoteHuitPoints} & \num{\frontazoteHuitProuves}
      & \num{\frontazoteHuitPlateau} palier, artefact de la chaîne d'amorçage \\
    \bottomrule
  \end{tabularx}
\end{table}

\begin{table}[H]
  \centering
  \caption{Réallocation entre le statu quo à enveloppe libre et le plafond de subventions le
  plus serré (\SI{\frontsubvQuatreunSeuil}{\euros}), en hectares par groupe \gls{rpg}.}
  \label{tab:ann-realloc}
  \small
  % Genere par memoire/build_chiffres.py -- ne pas editer.
\begin{tabular}{@{}lrrr@{}}
  \toprule
  \textbf{Groupe} & \textbf{Avant (ha)} & \textbf{Après (ha)} & \textbf{Écart (ha)} \\
  \midrule
  Prairie & \num{6096} & \num{7635} & \num{1539} \\
  Maraîchage & \num{1145} & \num{2237} & \num{1092} \\
  Agrumes & \num{9} & \num{657} & \num{649} \\
  Igname & \num{279} & \num{576} & \num{298} \\
  Jachère & \num{443} & \num{570} & \num{127} \\
  Vergers & \num{6} & \num{64} & \num{58} \\
  Ananas & \num{459} & \num{404} & \num{-55} \\
  Banane export & \num{2004} & \num{95} & \num{-1909} \\
  Canne à sucre & \num{12773} & \num{10541} & \num{-2231} \\
  \emph{autres groupes} & --- & --- & \num{-12} \\
  \bottomrule
\end{tabular}

\end{table}

\begin{table}[H]
  \centering
  \caption{Front marge $\times$ azote sous la configuration de référence, point par point.
  Pente moyenne \num{\frontazoteRefPente}~\eur/kg~N, dual local au point le plus serré
  \num{\frontazoteRefDualBas}~\eur/kg~N.}
  \label{tab:ann-front-azote-ref}
  \small
  % Genere par memoire/build_chiffres.py -- ne pas editer.
\begin{tabular}{@{}lrr@{}}
  \toprule
  \textbf{Plafond} & \textbf{Niveau atteint} & \textbf{Objectif (M\eur)} \\
  \midrule
  \num{1061997} & \num{1061851} & \num{69.12} \\
  \num{1158542} & \num{1158420} & \num{70.73} \\
  \num{1351632} & \num{1350263} & \num{74.12} \\
  \num{1544722} & \num{1542762} & \num{77.16} \\
  \num{1737813} & \num{1735470} & \num{79.75} \\
  \bottomrule
\end{tabular}

\end{table}

\begin{table}[H]
  \centering
  \caption{Front marge $\times$ azote sous la politique agroécologique P8, point par point.
  Pente moyenne \num{\frontazoteHuitPente}~\eur/kg~N, dual local au point le plus serré
  \num{\frontazoteHuitDualBas}~\eur/kg~N.}
  \label{tab:ann-front-azote-huit}
  \small
  % Genere par memoire/build_chiffres.py -- ne pas editer.
\begin{tabular}{@{}lrr@{}}
  \toprule
  \textbf{Plafond} & \textbf{Niveau atteint} & \textbf{Objectif (M\eur)} \\
  \midrule
  \num{810979} & \num{807152} & \num{63.30} \\
  \num{946142} & \num{943807} & \num{65.78}\dag \\
  \num{1081306} & \num{1066127} & \num{66.55}\dag \\
  \num{1216469} & \num{1066127} & \num{66.55}\dag \\
  \num{1351632} & \num{1351311} & \num{69.09}\textsuperscript{*} \\
  \bottomrule
\end{tabular}

\end{table}

\section{Coûts marginaux et prix duals}

\begin{table}[H]
  \centering
  \caption{Prix duals des contraintes nommées, lus sur la relaxation linéaire du run indiqué,
  et les trois réserves qui conditionnent leur emploi.}
  \label{tab:ann-duals}
  \scriptsize
  \setlength{\tabcolsep}{4pt}
  \begin{tabularx}{\textwidth}{@{}>{\raggedright\arraybackslash}p{34mm}X@{}}
    \toprule
    \textbf{Grandeur} & \textbf{Valeur ou réserve} \\
    \midrule
    Un point d'\glsentryshort{ift}, sous P7 & \SI{\dualIftSept}{\euros} \\
    Un point d'\glsentryshort{ift}, sous P8 & \SI{\dualIftHuit}{\euros} \\
    Un kilogramme d'azote, sous P8 & \SI{\dualAzoteHuit}{\euros} \\
    \midrule
    Un dual n'existe que sous une politique & les deux premières lignes sont la même grandeur
      physique sous deux politiques, et le prix va du simple à moins du double. Un prix dual
      cité sans le run dont il est issu ne veut rien dire \\
    Ce sont les duals de la relaxation linéaire & figer les variables entières ne fonctionne
      pas ici, le modèle étant purement binaire. La relaxation obtenue n'a plus aucune variable
      libre et tous les duals y valent zéro \\
    Un dual s'oppose à la pente du front correspondant & un écart important signale que le dual
      est lu hors de son voisinage de validité. C'est ce qui se produit sur l'enveloppe
      publique, où le dual au point le plus serré
      (\num{\frontsubvQuatreDualBas}~\eur/\eur) vaut près du double de la pente moyenne
      (\num{\frontsubvQuatrePente}~\eur/\eur). L'extrapoler sur tout l'intervalle surestimerait
      l'effet d'environ trois quarts \\
    \bottomrule
  \end{tabularx}
\end{table}
